\section{Consequences for numerical and measurement choices}
\label{sec:consequences}
The corrections change identifiable quantitative prescriptions, even where a
qualitative demonstration of successful trajectory reconstruction remains
valid. They should be assessed separately from the general possibility of
simulating a conservative density.

\paragraph{The geometric acceptance test.}
Under the printed definition, changing the packet width or taking a smaller
time step cannot make the stated directional-map ratio less than about $0.1$:
its lower bound is one. For the concrete nonlinear shear with
$\sigma=0.01$ and $h=0.001$, the checked raw ratio is $1.0000098691$, while the
increment ratio is $0.0044399602$. Thus the same step is rejected by the
printed statistic and can have small actual displacement distortion. A new
definition requires a newly justified interpretation of the threshold; merely
changing the right-hand numerical tolerance does not repair the missing
initial offset in the source's expansion. The complete-flow estimate above
states the additional temporal contribution that must be retained when
$\sigma$ is proportional to $h$.

\paragraph{The readout target and stopping precision.}
For the symmetric finite-grid packet with width $0.01$, its coordinate
variance is approximately $10^{-4}$ and its fringe contrast is approximately
$0.99980002$, but its phase maximizer remains the arithmetic mean. There is
no nonzero phase-bias floor in this example. A user should not stop refining
a mean estimate solely because the packet has nonzero variance. Conversely,
the explicit same-fringe pair has a mean difference of about
$1.966\times10^{-4}$ that no number of samples of that fixed-frequency fringe
can distinguish. One must change the encoding or obtain additional
information to estimate an arbitrary arithmetic mean. The source's direct
observable encoding is one available alternative. Whether more accurate
readout improves a \emph{single-trajectory} estimate still depends on the
separate cloud-to-trajectory bias.

\paragraph{Shot allocation.}
For independent samples with root-mean-square error target $\epsilon=0.001$,
a coordinate variance of $10^{-4}$ gives $M\approx100$, while the corresponding
ancilla variance $1/4$ gives $M=250{,}000$. The factor is $2500$, not a change
of the inverse-square exponent. These are variance-based root-mean-square
counts, not fixed-confidence guarantees. A confidence requirement introduces
its own tail bound; it does not identify the two random variables.

\paragraph{Packet-width calibration.}
Under the same source heuristic $R\lesssim0.1$, the printed decimal
$R/\sigma=2.29$ suggests $\sigma\lesssim0.04367$, whereas the displayed
definitions give $\sigma\lesssim0.14142$. This is a factor of about $3.24$
in the nominal width ceiling. Other localization, boundary, and error-budget
constraints may be stricter. Both calibrations accept the source's small
initial Lotka--Volterra width, so the correction is to the quantitative
guideline rather than a claim that this particular plotted orbit fails.

\paragraph{Observable error budgets and covariance constants.}
The $C^2$ counterexample has a cubic expectation in the packet width. A
fourth-order remainder budget underestimates it by an unbounded factor as the
width tends to zero; either the regularity or the budget must change.
Coordinate and quadratic moments remain exact identities. Appendix~
\ref{sec:covariance} identifies the precise smaller covariance bound that
fails and evaluates the repaired main-text constant in the same example.
